\begin{abstract}
Singular-value spectra are a compact description of how transformer weight matrices amplify and suppress internal directions, but their causal role in downstream behavior is still poorly understood. We study post-hoc spectral smoothing and spiking interventions on large language models by modifying only the singular values of MLP \texttt{up\_proj} and \texttt{down\_proj} matrices while keeping singular vectors fixed. The intervention preserves the geometric mean of singular values and continuously controls spectral inequality: positive perturbations sharpen dominant modes, whereas negative perturbations smooth the spectrum toward lower inequality. Across five Llama and Qwen instruction models, we evaluate the resulting systems on a fine-grained capability suite, auxiliary reasoning benchmarks, judge-scored truthfulness/hallucination/safety probes, and TRAIT personality-style measurements. To interpret subtask-level score shifts, we fit perturbation responses in a three-dimensional capability basis consisting of factual knowledge, language understanding, and deductive reasoning. The effects are structured rather than random: spectral perturbations induce model- and module-specific capability reweighting, selective task gains, coordinated TRAIT shifts, and layer-dependent sensitivity linked to pre-existing spectral inequality. These results position singular-value perturbation as a lightweight diagnostic intervention for probing how capabilities and behavioral tendencies are organized in pretrained LLM weights.
\end{abstract}

\section{Introduction}

Large language model behavior is often analyzed through prompts, finetuning, activation interventions, or output-level benchmark scores. These views are indispensable, but they leave open a complementary question: can a simple, interpretable perturbation of pretrained weight geometry reveal how model capabilities are organized? Transformer MLP projections are natural targets for this question. If \(W = U\Sigma V^\top\), each singular value \(\sigma_i\) determines the gain with which the input direction \(v_i\) is mapped into the output direction \(u_i\). The singular-value spectrum therefore describes not only the effective rank of a matrix, but also the hierarchy of internal directions that a layer amplifies or suppresses.

We use this observation to study a controlled weight-space intervention. For every MLP \texttt{up\_proj} and \texttt{down\_proj} matrix, we keep the singular vectors fixed and perturb only the log singular values around their mean:
\[
\ell_i^{(\alpha)}=\bar{\ell}+(1+\alpha)(\ell_i-\bar{\ell}),
\qquad
\ell_i=\log\sigma_i,\qquad
\bar{\ell}=\frac{1}{r}\sum_i \ell_i .
\]
Positive \(\alpha\) amplifies deviations from the mean log singular value, producing a more peaked spectrum; negative \(\alpha\) shrinks the spectrum toward its geometric mean, producing a smoother spectrum. Because \(\bar{\ell}\) is preserved, the geometric mean of singular values remains unchanged. The intervention therefore changes spectral inequality without simply rescaling the matrix and without changing the singular-vector basis.

The goal is not to find a universally beneficial perturbation strength. Instead, we ask whether spectral smoothing and spiking expose structured behavioral tradeoffs. A priori, there are several plausible outcomes. If the singular-value distribution mostly reflects redundant numerical degrees of freedom, perturbing it should mainly damage performance. If high-gain and low-gain directions support different computational roles, then perturbing spectral inequality should produce non-uniform changes across tasks, capabilities, and models. The latter view is motivated by recent interpretability and spectral analyses showing that singular directions can correspond to meaningful semantic or functional subspaces, while smaller singular directions may still encode adaptation, alignment, routing, or other low-visibility signals.

We evaluate this intervention across five open instruction models: Llama 3.2 1B, Llama 3.2 3B, Llama 3.1 8B, Qwen3 8B, and Qwen3 30B-A3B MoE. The main analysis uses a fine-grained capability suite consisting of 118 subtasks from MMLU-Pro, MMLU-Redux, AGIEval English, and BBH. Each subtask is assigned a judge-derived capability-demand mixture. A seven-dimensional rubric is first used for scoring task demands, and a PCA/simplex analysis shows that the benchmark-wide geometry is well summarized by three exposed axes: factual knowledge, language understanding, and deductive reasoning. We then fit score changes relative to \(\alpha=0\) in this three-dimensional capability basis.

The resulting picture is not a uniform quality loss. Spectral perturbations reconfigure behavior across benchmark modules, auxiliary tasks, and personality-style probes. Large positive perturbations more reliably suppress language and deductive responses, whereas moderate negative perturbations can preserve or selectively improve factual or task-local behavior. Auxiliary benchmarks show narrow gains on some tasks, with diagnostic error decomposition indicating that several no-CoT improvements correspond to reduced wrong-answer rates rather than parser artifacts. TRAIT probes also move in coordinated ways with fitted capability responses, suggesting that benchmark accuracy shifts and personality-like behavioral summaries are coupled under the same spectral intervention.

Layer-wise spectral analysis explains why a single global \(\alpha\) can have heterogeneous effects. Different layers and projections start from different degrees of spectral inequality, so the same nominal perturbation induces different realized changes in top singular values and Gini coefficients. In MoE models, expert-level variation adds another source of non-uniformity. Thus, spectral smoothing and spiking should be interpreted as structure-dependent interventions rather than scalar noise injections.

Our contributions are:
\begin{itemize}
    \item We introduce an energy-preserving singular-value smoothing/spiking intervention for transformer MLP projections that modifies spectral inequality while keeping singular vectors fixed.
    \item We evaluate the intervention across five Llama and Qwen instruction models using a fine-grained capability suite, auxiliary reasoning tasks, judge-scored behavioral probes, TRAIT measurements, and output-distribution diagnostics.
    \item We fit subtask responses in a three-dimensional capability basis and show that spectral perturbations induce model- and module-specific capability reweighting rather than uniform degradation.
    \item We connect behavioral heterogeneity to layer-wise spectral non-uniformity, showing that layers with different initial singular-value distributions receive different realized update magnitudes under the same global \(\alpha\).
\end{itemize}

\section{Related Work}

\paragraph{Singular directions as interpretable transformer structure.}
Several recent works suggest that singular directions of transformer weight matrices are not merely numerical artifacts. The singular-value decompositions of transformer MLP and attention matrices have been shown to expose interpretable semantic directions, especially among high-gain directions in later layers~\citep{svd_transformer_interpretable}. More recent circuit-level work decomposes attention and MLP components into singular-vector subspaces and shows that functional roles can be distributed across compact directions inside components that would otherwise be treated as atomic~\citep{beyond_components_svd_circuits}. These studies motivate our use of SVD as an intervention basis. We differ by perturbing singular values directly and measuring downstream behavioral changes, rather than only interpreting existing directions.

\paragraph{High- and low-spectrum components.}
Prior work does not support a simple mapping such as ``high singular values encode logic'' and ``low singular values encode knowledge.'' A more defensible view is that large singular directions often correspond to dominant, strongly expressed semantic or computational modes, while smaller singular directions can still carry task-relevant, alignment-relevant, or low-visibility information. Random-matrix analyses show that small singular values in trained transformers can deviate meaningfully from noise baselines and may become especially important after finetuning or alignment~\citep{small_singular_values_matter}. Related spectral work identifies low-spectrum ``dark'' signals that are weakly visible to direct decoding yet functionally relevant to attention sinks and global routing behavior~\citep{spectral_filters_dark_signals}. Conversely, layer-selective rank reduction can sometimes improve reasoning-like benchmarks by retaining dominant low-rank structure and suppressing competing or noisy components~\citep{laser}. Our smoothing/spiking perturbation provides a continuous alternative to hard truncation: it preserves all directions but changes their relative gains.

\paragraph{Singular-value adaptation and low-rank finetuning.}
Parameter-efficient finetuning methods also point to differentiated roles across spectral bands. PiSSA initializes adaptation using principal singular values and vectors, treating high-spectrum directions as strong carriers of pretrained structure~\citep{pissa}. MiLoRA instead targets minor singular components, motivating the low-spectrum subspace as a way to adapt long-tail or task-specific behavior while reducing interference with the pretrained backbone~\citep{milora}. KaSA frames singular-value adaptation as a mechanism for activating task-relevant parametric knowledge~\citep{kasa}. These works use spectral information for training-time adaptation, whereas our intervention is post-hoc and diagnostic: no adapter is trained, and the same pretrained model is evaluated under controlled spectral reweighting.

\paragraph{Parametric knowledge in MLPs.}
A separate line of interpretability work locates factual knowledge and factual editing mechanisms in transformer feed-forward layers. Feed-forward networks have been interpreted as key-value memories whose keys detect textual patterns and whose values affect output-token distributions~\citep{geva_kv_memory}. ROME uses causal tracing and rank-one editing to localize and modify factual associations in GPT-style models~\citep{rome}, while MEMIT extends this view to mass editing~\citep{memit}. Knowledge-neuron work similarly identifies neurons associated with factual recall~\citep{knowledge_neurons}. These studies support the relevance of MLP weights to parametric knowledge, but they do not specify which singular-value bands carry particular knowledge types. We therefore treat factual knowledge as an empirically fitted capability response rather than assuming it occupies either the top or tail of the spectrum.

\paragraph{Capability decomposition beyond aggregate scores.}
Aggregate benchmark scores can obscure heterogeneous subtask behavior. A perturbation may improve one subtask family while harming another, and equal score changes can reflect different mixtures of factual recall, language comprehension, and reasoning. Our fine-grained capability suite follows this motivation by assigning each subtask a judge-derived capability-demand mixture and fitting perturbation-induced score deltas in a low-dimensional capability basis. This analysis is complementary to standard accuracy reporting: it allows us to distinguish broad capability shifts from localized module stress and to connect task-level behavior with TRAIT and judge-scored probes.

\paragraph{Behavioral probes and trait-style evaluation.}
Truthfulness, hallucination, safety, and personality-style probes are often used to characterize alignment-relevant or stylistic behavior beyond raw benchmark accuracy. In this paper, TRAIT scores are not interpreted as stable human-like personalities. Instead, they serve as compact behavioral summaries that can be compared across perturbation strengths. By analyzing TRAIT shifts together with fitted capability responses, we test whether spectral interventions move benchmark performance and broader behavioral tendencies along coordinated axes. This connects weight-space spectral structure to a wider behavioral profile than task accuracy alone.
